\begin{python0}
from solutions import *; clear()
\end{python0}
\sectionthree{Summary}

A \textbf{statement} is something that will cause your computer to perform some operation(s).

In C++ statements must terminate with a semi-colon.

C++ is case-sensitive. 

C++ ignores whitespace (between tokens).

Something that looks like \verb!"Hello, World!\n"! (i.e. characters within double quotes) is called a \textbf{string} or a \textbf{C-string}.

A \textbf{string} is either empty (the null string), i.e. \verb!""!, or it is made up of characters. 

For instance the first character of the string
\begin{console}[frame=none]
 
    "Hello, World!\n"

\end{console}
is \verb!'H'!.

There are special characters: 
\verb!'\n'! causes the cursor to jump to a new line while 
\verb!'\t'! moves to the next tab position, 
\verb!'\"'! is the double-quote character, 
and \verb!'\''! is the single-quote character. 

The backslash character is \verb!'\\'!. These are called \textbf{escape characters}.

To print a string, you execute the following statement:
\begin{console}[frame=none, commandchars=\~\!\@]

        std::cout << ~textit![some string]@;

\end{console}
You can print more than one string in a single print statement:
\begin{console}[frame=none, commandchars=\~\!\@]

        std::cout << ~textit![string1]@ << ~textit![string2]@
                  << ~textit![string3]@ << ~textit![string4]@;

\end{console}
Printing one or more characters is the same:
\begin{console}[frame=none, commandchars=\\\{\}]

        std::cout << \textit{[some character]};
        std::cout << \textit{[character1]} << \textit{[character2]}
                  << \textit{[character3]} << \textit{[character4]};

\end{console}

You can mix print strings and characters. For instance
\begin{console}[frame=none, commandchars=\\\{\}]

        std::cout << \textit{[some character]}
	          << \textit{[some string]};

\end{console}
or
\begin{console}[frame=none, commandchars=\\\{\}]

        std::cout << \textit{[some string]}
	          << \textit{[some character]};

\end{console}
You can force a newline by printing \verb!std::endl!:
\begin{verbatim}
        std::cout << std::endl;
\end{verbatim}



If your program has two or more statements like this:
\begin{console}[frame=none, commandchars=\\\{\}]

        \textit{[statement 1]};
        \textit{[statement 2]};
        \textit{[statement 3]};

\end{console}
then C++ executes top to bottom (for now).

